\subsection{Cel dokumentu}

Każdy rozdział znajduje się w osobnym katalogu, a jego plik \texttt{main.tex}
dołącza kolejne podrozdziały poleceniem \verb|\input|. Dzięki temu pliki są krótkie,
a zmiany w historii Git łatwe do przejrzenia.

Odwołania do literatury wstawia się poleceniem \verb|\cite|, \np{}~\cite{knuth1984texbook}.
Przykłady rysunków, tabel i wzorów znajdują się w \cref{sec:przyklady}.
